\section{A first example: the integers under addition}\label{sec:06-groups:03-integers-example:1}

Consider packaging \((\mathbb{Z}, +, 0, -(-))\) as a \texttt{Group Int}. The construction proceeds field by field, proving each obligation with a short, explicit tactic proof.

\begin{lstlisting}
def intGroup : Group Int where
  op := fun a b => a + b
  id := 0
  inv := fun a => -a
  assoc := by
    intro a b c
    -- Goal: (a + b) + c = a + (b + c)
    exact Int.add_assoc a b c
  id_left := by
    intro a
    -- Goal: 0 + a = a
    exact Int.zero_add a
  id_right := by
    intro a
    -- Goal: a + 0 = a
    exact Int.add_zero a
  inv_left := by
    intro a
    -- Goal: (-a) + a = 0
    exact Int.add_left_neg a
  inv_right := by
    intro a
    -- Goal: a + (-a) = 0
    exact Int.add_right_neg a
\end{lstlisting}

Each field is proved separately. Each proof is a single \texttt{intro} (to name the universally quantified variables) followed by \texttt{exact} naming the exact core-library lemma that already states this fact about \texttt{Int}. No step is hidden: \texttt{Int.add\_assoc}, \texttt{Int.zero\_add}, and the rest are themselves proved (elsewhere, in Lean's core library) by induction on the same \texttt{Nat}/\texttt{Int} representation introduced in Chapter 4. These lemmas are reused here rather than re-deriving integer arithmetic from scratch.

\begin{mathreading}

This shows \((\mathbb{Z}, +, 0, -)\) as an object of the category \(\mathbf{Grp}\). The term \texttt{intGroup} is a \emph{proof that \(\mathbb{Z}\) is a group}, built by giving the data \((+, 0, -)\) and verifying each axiom. The verification is not re-proved here but \emph{cited}: associativity is \(\mathrm{Int.add\_assoc}\), the identity laws are \(0 + a = a
= a + 0\), and the inverse laws are \((-a) + a = 0 = a + (-a)\). In textbook terms this is the one-line remark ``\(\mathbb{Z}\) under addition is an abelian group,'' with the underlying lemmas about \(\mathbb{Z}\) (themselves ultimately inductions on the integers) written out in full instead of just assumed.

\end{mathreading}

\textbf{Mathlib equivalent.} Mathlib requires no \texttt{intGroup}-style bundle at all: \texttt{Int} is \emph{already} registered as an \href{https://loogle.lean-lang.org/?q=AddCommGroup}{\texttt{AddCommGroup}} instance, and the five axioms above are available as free-standing lemmas that apply to every additive group, not just \texttt{Int}:

\begin{lstlisting}
example : AddCommGroup Int := inferInstance

example (a b c : Int) : (a + b) + c = a + (b + c) := add_assoc a b c
example (a : Int) : 0 + a = a := zero_add a
example (a : Int) : a + 0 = a := add_zero a
example (a : Int) : -a + a = 0 := neg_add_cancel a
example (a : Int) : a + -a = 0 := add_neg_cancel a
\end{lstlisting}

This is the same content as \texttt{intGroup}, the same five facts about \(\mathbb{Z}\). But where the book \emph{assembles} a \texttt{Group Int} term by hand, Mathlib's version has nothing to assemble: the instance already exists, found automatically by \href{https://loogle.lean-lang.org/?q=inferInstance}{\texttt{inferInstance}}. And \href{https://loogle.lean-lang.org/?q=add_assoc}{\texttt{add\_assoc}}/\href{https://loogle.lean-lang.org/?q=zero_add}{\texttt{zero\_add}}/etc. are generic lemmas about \emph{any} \texttt{AddCommGroup}, so they read somewhat differently from \texttt{Int.add\_assoc}: they apply equally well to Chapter 6's \texttt{perm3Group}-style examples once those are phrased in Mathlib's \texttt{Group}/\texttt{AddCommGroup} classes (Chapter 6 §4 does exactly that next).
